% Felipe Muñoz Casasola
% This is free and unencumbered software released into the public domain.

% Anyone is free to copy, modify, publish, use, compile, sell, or
% distribute this software, either in source code form or as a compiled
% binary, for any purpose, commercial or non-commercial, and by any
% means.

% In jurisdictions that recognize copyright laws, the author or authors
% of this software dedicate any and all copyright interest in the
% software to the public domain. We make this dedication for the benefit
% of the public at large and to the detriment of our heirs and
% successors. We intend this dedication to be an overt act of
% relinquishment in perpetuity of all present and future rights to this
% software under copyright law.

% THE SOFTWARE IS PROVIDED "AS IS", WITHOUT WARRANTY OF ANY KIND,
% EXPRESS OR IMPLIED, INCLUDING BUT NOT LIMITED TO THE WARRANTIES OF
% MERCHANTABILITY, FITNESS FOR A PARTICULAR PURPOSE AND NONINFRINGEMENT.
% IN NO EVENT SHALL THE AUTHORS BE LIABLE FOR ANY CLAIM, DAMAGES OR
% OTHER LIABILITY, WHETHER IN AN ACTION OF CONTRACT, TORT OR OTHERWISE,
% ARISING FROM, OUT OF OR IN CONNECTION WITH THE SOFTWARE OR THE USE OR
% OTHER DEALINGS IN THE SOFTWARE.

% For more information, please refer to <https://unlicense.org/>

\documentclass{article}

\usepackage[spanish]{babel}
\usepackage[margin=3cm]{geometry}
\usepackage{microtype}
\usepackage[pdfusetitle]{hyperref}
\usepackage[utf8]{inputenc}
\usepackage{minted}
\usepackage{verbatim}

\title{Geany}
\author{Felipe Muñoz Casasola}
\date{21 de marzo de 2026}

\begin{document}
    \maketitle
    \clearpage
    
    \section{Instalación}
    Lo mejor es compilarlo desde 0. Para ello puedo usar toolbox para
    tener ahí dentro un contenedor de mi ordenador y compilarlo
    instalando los paquetes necesarios sin tener que instalármelos en
    mi sistema anfitrión para compilarlo. Para más información sobre
    eso puedo mirar la guía \texttt{Perfiles-de-trabajo.pdf}.

    \section{Complementos}
    Es posible que si lo hago desde un contenedor haya algún problema
    con \texttt{pkg-config}. Para lo que tengo que copiar el archivo
    \texttt{geany.pc} que hay en la carpeta donde está el código
    fuente de Geany, y actualmente lo tengo que copiar en
    \texttt{/usr/lib/pkgconfig/}.

    También tengo que comprobar que estén activados los Complementos
    que quiero. Y si no fuera así ejecutar
    \texttt{./configure --enable-complemento} para ver por qué no
    está activado e instalar las dependencias que faltan (con
    \texttt{apt search}).

    \section{Configuración}
    Actualmente el único complemento que tengo es «Project organizer».
    Tengo el siguiente añadido en mi archivo
    \texttt{\$RUBY/.config/geany/snippets.conf}:

    \inputminted[breaklines, breakanywhere, breaksymbolleft="", breaksymbolright="", breakanywheresymbolpre="", breakanywheresymbolpost=""]
    {systemd}{inc/snippets.conf}

    El tema del editor es Dracula y mi archivo de configuración tiene el
    siguiente contenido:

    \inputminted[breaklines, breakanywhere, breaksymbolleft="", breaksymbolright="", breakanywheresymbolpre="", breakanywheresymbolpost=""]
    {systemd}{inc/geany.conf}

    Y en la wiki de Geany hay guías, en concreto por el momento
    solamente hay una para tener un analizador estático de código
    en Ruby. Pero quizá añadan más, está en \href{https://wiki.geany.org/start}
    {https://wiki.geany.org/start}.

    \section{Rubocop}
    Esta aplicación instalada desde \texttt{sudo gem install rubocop} hace
    una comprobación del estilo y la sintaxis en Ruby. Tal y como lo
    tengo configurado, arregla automáticamente los errores de estilo
    que no afectan en ningún caso a la lógica del programa. La
    configuración se encuentra en
    \texttt{\$RUBY/.config/rubocop/rubocop.conf}.
    
    \inputminted[breaklines, breakanywhere, breaksymbolleft="", breaksymbolright="", breakanywheresymbolpre="", breakanywheresymbolpost=""]
    {yaml}{inc/rubocop.yml}
\end{document}
